\section{Limitations}\label{sec:limitations}

\paragraph{Gating and family coverage.}
The census covers only what is publicly downloadable without license acceptance.
Gating excluded \kvcGatedSkipped{} repositories, and the exclusions are
concentrated: every \texttt{meta-llama} checkpoint was gated at snapshot time,
so the census contains no rows from that organization, and the earlier
\texttt{google/gemma} catalog was skipped, leaving only the recent
\texttt{gemma-4} checkpoints to represent that family. Llama-architecture checkpoints released by other included publishers are present, but the primary Llama releases are absent. Both families anchor a
large share of real deployments, and the Gemma~3 report documents exactly the
window-heavy, KV-saving designs this paper studies \cite{gemmateam2025gemma3}.
Their partial absence shifts yearly medians and the mechanism mix in directions
we cannot sign from the data at hand; a Llama-only or Gemma-only census would
likely sit at different medians for both KV per token and window share.

\paragraph{Release-date proxy and snapshot.}
Release dating uses the Hub \texttt{createdAt} timestamp, which measures upload,
not publication. Re-uploads, staged rollouts, and repository reorganization can
move a model across year boundaries, and every per-year statistic inherits that
noise. A model announced in one year and uploaded in the next is dated late;
variants uploaded together with their base model are dated early. The snapshot is single-point-in-time (2026-10-10): repositories were not
re-crawled, deleted or renamed repos are not re-checked, and 2026 figures cover
a partial year of \kvcRqOneNTwentySix{} deduplicated rows, so the 2026 medians
will move once the year completes.

\paragraph{Dropped rows and dedup heuristics.}
Coverage inside included repositories is imperfect. \kvcKvErrRows{}
configurations could not be parsed into an attention geometry (missing or
non-standard fields) and were dropped. \kvcNoParams{} rows lack parameter counts
and are excluded from per-parameter statistics and from every capacity
computation. The dedup key is organization, name stem, and geometry: it collapses base/instruct twins
but \emph{not} same-geometry fine-tunes with distinct names, so the
\kvcDedupRows{} pre-registered groups realize only \kvcGeoRows{} distinct
geometries (\kvcWorstGeoCluster{} rows share the most common one). Trend medians,
p-values, and two verdicts (H1, H4) are sensitive to this unit choice, which is why
every affected statistic is reported under both units. Multimodal rows (\kvcMMRows{}
deduplicated rows) carry vision-tower-inflated parameter counts against
language-backbone KV, deflating their per-parameter cost; RQ1 reports the
multimodal-excluded cut for this reason.

\paragraph{Analytic values are theoretical minima.}
All KV figures follow from the configurations, not from measurements of deployed
runtimes. Implementations can allocate more: latent attention may be cached
decompressed, sliding windows may be allocated in full, and recurrent states
carry padding. RQ4 quantifies the gap on reduced-depth instances:
\kvcValWithin{} of \kvcValRows{} validation rows agree within 1\%, spanning
\kvcValClasses{} attention classes, and the largest residual comes from a pure
recurrent model whose state accounting differs (Table~\ref{tab:validation}).
The validation covers classes, not every census family.

\paragraph{Capacity model idealizations.}
The capacity model is deliberately minimal. It ignores allocator fragmentation
--- the problem paged attention was built to solve \cite{kwon2023paged} ---
plus activation memory, kernel workspaces, and other runtime internals. Runtime
overhead enters as a modeled parameter reported as a sensitivity grid (0--12
GiB; headline tables use 6 GiB), not as a measurement. Budgets are treated as
uniform unified memory; prefill/decode disaggregation changes the arithmetic
\cite{zhong2024distserve}. Finally, the Hub is a biased frame. We include original
publishers only (\kvcOrgCount{} organizations), open weights only, and no closed
models; the fine-tune and merge ecosystem is excluded by design, and recent
release years skew toward many small checkpoints. Fleet medians therefore
describe the open-weight release stream, not production traffic, and readers
should not read them as market shares.
